% RESEARCH_DRAFT_SHA256: e645fc5c2f2c4ec760ab2dc656d6f065b016156ac1681cbf85ef98c316b0cbf7
% Layout and source-reference wrappers expanded for the web renderer.
\documentclass[11pt]{article}
\usepackage[margin=25mm]{geometry}
\usepackage{amsthm}
\usepackage{fontspec,unicode-math}
\setmainfont{Latin Modern Roman}
\setmathfont{Latin Modern Math}
\usepackage{amsmath,amssymb,mathtools,xurl,hyperref}
\hypersetup{colorlinks=true,urlcolor=blue,linkcolor=blue}
\setcounter{secnumdepth}{0}
\setlength{\parindent}{0pt}
\setlength{\parskip}{5pt}
\setlength{\emergencystretch}{3em}
\newcommand{\sref}[2]{\hyperlink{source:#1:#2}{[S#2]}}
\newcommand{\cataloguescope}[1]{\paragraph{Source record.}#1}
\newcommand{\floor}[1]{\left\lfloor#1\right\rfloor}
\newtheorem{proposition}{Proposition}
\newtheorem{lemma}[proposition]{Lemma}
\begin{document}
\paragraph{Research manuscript.} OpenAI \texttt{gpt-6-astra}, reasoning \texttt{ultra}; 11 October 2026.
\paragraph{Preserved source body.} The complete catalogue mathematical text and sources follow with layout and source links expanded, including source qualifications and its historical preparation notes. New research begins at the section headed Research attempt.

% UnsolvedMath ID 20002140; source catalogue code AIM-INFRASTRUCTURE-0040
\section{A replicable and inclusive developer on-ramp}\label{20002140}
{\small UnsolvedMath ID 20002140 · Computer Science}
\subsection{Definitions and mathematical statement}

\paragraph{Shared notation.}Standard number systems are $\mathbb N,\mathbb Z,\mathbb Q,\mathbb R,\mathbb C$. Specialized notation and index ranges follow the cited source and the definitions below.


An on-ramp is an introductory path that enables a beginner to complete a specified procedure. Prerequisites describe required prior knowledge; a reproducible instruction sequence identifies tools, steps, and expected results. Inclusive design considers barriers experienced by the intended users rather than assuming identical backgrounds.
\paragraph{Record scope and qualifications.}
The catalogue title is retained as metadata. Named certificates, contracts, bounds or frameworks in that title are not asserted as results of the original source extract.
The registry marks this wording reconstructed. The software, procedure, beginner population, and usability criteria are not supplied.

Catalogue screening limitation: Workshop planning or software infrastructure agenda; the supplied passage does not identify a determinate mathematical target.

The source snapshot labels this record partially solved; this is catalogue metadata rather than independent verification of a mathematical result or the current completion of a project.
The source text below is rendered literally where necessary; damaged formulas and composite extracts are not silently repaired into a new mathematical claim.
The supplied catalogue marks this formulation as an unverified reconstruction; missing assumptions are not certified by the source link.
\paragraph{Statement recorded in the supplied source.}
The goal of this group is to produce an example of an "on-ramp" which could be given to someone with little to no experience in software, from which they could learn to do some procedure. The guiding principles of this group is that the output should be considerate of people from a marginalized background, and should aim to create something that would be replicable. In addition, the group, to the extent that they can, should attempt to finish a complete document. \href{http://aimpl.org/cyberinfrastructure/3/}{[S2]}
\paragraph{Source status.}
The supplied catalogue labels this source record partially solved; that metadata was checked against the supplied snapshot on 2026-10-11. The label is preserved for provenance and is not an independent verification of present open status or mathematical correctness.
\subsection{Short English statement}
Produce an accessible and reproducible introduction to a software procedure for beginners.
\subsection{Sources}
\begin{itemize}
\item[S1] Ulam AI and the UnsolvedMath Contributors, UnsolvedMath: A Curated Collection of Open Mathematics Problems, record 20002140 (AIM-INFRASTRUCTURE-0040). Catalogue curation and metadata: CC BY 4.0. The source extract is retained; definitions, English summaries and qualifications are editorial additions.
\url{https://huggingface.co/datasets/ulamai/UnsolvedMath}.
\item[S2] AIM Workshop Problem Lists, source linked by catalogue record AIM-INFRASTRUCTURE-0040.
\url{http://aimpl.org/cyberinfrastructure/3/}.
\end{itemize}
\paragraph{Preparation and review.}
The original source and UnsolvedMath contributors retain attribution for the source extract. Codex (GPT-6; exact serving revision not exposed) assisted with definitions, English summaries, scope notes and formatting. No new human mathematical authorship or independently verified solution is asserted. Permanent source ID 20002140 and code AIM-INFRASTRUCTURE-0040 checked against the complete supplied registry and current repository. Distinct upstream records, including repeated questions and broad or damaged extracts, are retained for complete coverage of the source CS category. A repeated record is not a new independent conjecture. Complete mathematical-equivalence and current-literature review has not been performed. Preparation and final review follow contribution prompts C and D, with source gaps disclosed. Structural and rendering checks do not establish a new theorem or open status.
\label{end:20002140}

\section{Research attempt}
\subsection{Literature, source context and scope}
The complete catalogue body above retains the original agenda, source status, and historical preparation. The AIMPL problem page could not be retrieved. The supplied complete source transcription was read, and the accessible primary AIM workshop announcement was checked for context. The announcement is not a substitute for verification of the inaccessible problem page. The catalogue's earlier proposed frameworks and research summaries are context, not established solutions credited to this attempt. The following is an explicitly limited investigation of one route; it neither substitutes a formalized toy model for the entire source agenda nor claims empirical or deployed-software results. Elementary auxiliary facts are not new original-target progress.

\subsection{A complete small procedure and its limited accessibility claim}
The source asks for an actual beginner on-ramp. Here is a self-contained paper-executable procedure for finding a greatest common divisor, rather than only an outline. The prerequisites are recognizing nonnegative whole numbers and dividing with a remainder. No account, installation, timed task, or spoken instruction is required. These design choices are explicit, but their accessibility has not been tested with users.
Choose two nonnegative whole numbers $a,b$, not both zero. Keep a two-column written record labeled $a$ and $b$. If $b=0$, stop and report $a$. Otherwise divide $a$ by $b$ to obtain the remainder $r$, the unique integer with $0\le r<b$ and $a=qb+r$ for an integer $q$. On a new line write the old $b$ in the first column and $r$ in the second, then repeat. If the input contains fractions or negative numbers, it is outside this procedure's stated input contract. If the remainder is not smaller than the divisor, redo that division before continuing.
\begin{proposition}
This procedure terminates and reports $\gcd(a,b)$ for every permitted input.
\end{proposition}
\begin{proof}
For integers $a=qb+r$, a positive integer divides both $a,b$ if and only if it divides both $b,r$: subtraction gives one direction and addition the other. Thus each new line preserves the set of common positive divisors. At a nonterminal step the second entry strictly decreases and remains nonnegative. A strictly decreasing sequence of nonnegative integers cannot continue indefinitely, so some line has second entry zero. Its first entry is positive, because a pair not both zero never becomes $(0,0)$ under the update. The common positive divisors of $(a,0)$ are precisely the positive divisors of $a$, whose greatest is $a$. The preserved divisor set proves correctness.
\end{proof}
For a worked run, $(252,105)$ becomes $(105,42)$, then $(42,21)$, then $(21,0)$, so the answer is 21. A reader can check the divisions $252=2\cdot105+42$, $105=2\cdot42+21$, and $42=2\cdot21$.
\subsection{Verification and remaining gaps}
The procedure was implemented in a short Python check and tested against the standard-library gcd on all 440 nonzero pairs with entries from 0 through 20. Mathematical correctness does not establish that the document is understandable, inclusive, or sufficient to onboard someone into software development. No participant evaluation or complete software-development on-ramp was conducted. The task is one concrete instructional example, deliberately smaller than the source's community goal.

\section{References and source attribution}
Primary source: \url{http://aimpl.org/cyberinfrastructure/3/}. Attempted primary AIMPL page retrieval; failed with timeout or gateway error. The supplied full source transcription was read; the inaccessible live page was not independently inspected.

\url{https://aimath.org/pastworkshops/cyberinfrastructure.html}. Read the primary December 2023 workshop announcement; it describes the open-source mathematical infrastructure agenda.

All retrievals were checked on 11 October 2026. Original contributor and catalogue attribution remains in the preserved Sources section. The elementary arguments below are self-contained; their exposition is not a claim of first discovery.


\paragraph{Integration review.} Codex (GPT-6; exact serving revision not exposed) reviewed the written arguments and expanded layout and source-reference wrappers. Repository metadata and the production web renderer were checked using Python 3.13 and JavaScript with MathJax 3.2.2. This model-assisted review is not independent human proof verification; no PDF compilation or formal proof checking is asserted.

\section{Completion Estimate}
\noindent\textbf{COMPLETION: 0\%}
No substantiated new progress toward the entire original target is claimed. The worked results and failed route are delimited above; a known result, illustrative model, or auxiliary obstruction is not counted as original-target completion.

\section{Final status and attribution}
\textbf{Attempt status: unresolved.} This describes this research attempt, not a new ruling on the source catalogue's status or on external literature claims.
The mathematics and this self-audit were generated by OpenAI \texttt{gpt-6-astra}, reasoning \texttt{ultra}, on 11 October 2026. Codex (GPT-6; exact serving revision not exposed) orchestrated the batch and contributed formatting/repository review. No human mathematical authorship, independent human proof checking, or formal verification is claimed. The historical preparation paragraphs in the reproduced catalogue body concern that body only.

Contribution prompt A was applied to the complete definition and supplied source record, followed by the adversarial checks of prompt B. Local reproduction files are \texttt{batch013/checks.py}, \texttt{computations.json}, and the per-ID audit JSON. They were not public links at preparation time. Numerical checks are finite and do not prove unrestricted results. No LaTeX engine was found locally; compilation was not performed by this research agent.

\end{document}
